\documentclass[12pt]{article}
\usepackage[utf8]{inputenc}
\usepackage[T1]{fontenc}
\usepackage[margin=1in]{geometry}
\usepackage{setspace}
\setlength{\parindent}{2em}
\setlength{\parskip}{1.25em}
\renewcommand{\baselinestretch}{1.0}
\makeatletter
\newcommand{\@subtitle}{}
\newcommand{\subtitle}[1]{\renewcommand{\@subtitle}{#1}}
\renewcommand{\@maketitle}{%
  \newpage
  \null
  \vskip 1em%
  \begin{center}%
    {\LARGE \@title \par}%
    \ifx\@subtitle\@empty\else
      \vskip 0.5em%
      {\large \@subtitle \par}%
    \fi
  \end{center}%
  \par
  \vskip 1em}
\makeatother
\usepackage{xcolor}
\usepackage[colorlinks=true,urlcolor=blue!55!black,linkcolor=black]{hyperref}

\title{\textit{Christian Legal Society Chapter of the University of California, Hastings College of the Law v. Martinez}, 561 U.S. 661 (2010)}
\subtitle{Case Brief}
\author{}
\date{}

\begin{document}
\maketitle

\section{Memory Jogger}

A public law school may condition official student-organization recognition and its attendant funding and communication benefits on an ``all-comers'' open-membership policy without violating the First Amendment, because the policy is a reasonable, viewpoint-neutral condition on access to a limited public forum.

\section{Detailed Case Facts}

The University of California, Hastings College of the Law (``Hastings''), founded in 1878 as the first law school in the University of California public-school system, operated a ``Registered Student Organization'' (``RSO'') program to encourage extracurricular student associations (\href{https://www.courtlistener.com/opinion/150544/christian-legal-soc-chapter-of-univ-of-cal-hastings-college-of-law-v/}{CourtListener}; \href{https://supreme.justia.com/cases/federal/us/561/661/}{Justia}). Official RSO status conferred substantial benefits: eligibility to apply for financial assistance drawn from a mandatory student-activity fee paid by all students, access to official Law School communication channels (a weekly Office of Student Services newsletter, designated bulletin boards, a Hastings-organization e-mail address, and participation in the annual Student Organizations Fair), permission to reserve Law School facilities for meetings and office space, and a license to use the Hastings name and logo.

In exchange for these benefits, RSOs had to be non-commercial organizations composed solely of Hastings students, submit their bylaws for approval, and comply with Hastings' written Nondiscrimination Policy (adopted in 1990), which tracked California's statutory ban on discrimination in public higher education (Cal. Educ. Code \S~66270) and prohibited unlawful discrimination on the basis of ``race, color, religion, national origin, ancestry, disability, age, sex or sexual orientation.'' Crucially, in a December 2005 Joint Stipulation of Facts submitted at summary judgment, both parties agreed that Hastings interpreted this policy in the RSO context to mandate acceptance of ``all comers'': RSOs had to ``allow any student to participate, become a member, or seek leadership positions in the organization, regardless of [her] status or beliefs,'' meaning, for example, that the Hastings Democratic Caucus could not bar Republican students from membership or leadership (\href{https://tile.loc.gov/storage-services/service/ll/usrep/usrep561/usrep561661/usrep561661.pdf}{U.S. Reports PDF}).

At the start of the 2004--2005 academic year, student leaders of a predecessor Christian RSO (the Hastings Christian Fellowship, which had operated as an RSO for a decade) affiliated with the national Christian Legal Society (``CLS-National''). National affiliation required the Hastings chapter (``CLS'') to adopt bylaws mandating that voting members and officers sign a ``Statement of Faith'' affirming orthodox Christian doctrines and conduct their lives accordingly. Following a 2004 national resolution, CLS interpreted its bylaws to exclude from voting membership and leadership anyone who held religious beliefs differing from the Statement of Faith or who engaged in ``unrepentant homosexual conduct'' (defined as sexual activity outside marriage between a man and a woman), though all students were welcome to attend meetings as guests.

On September 17, 2004, CLS applied for RSO status. Judy Hansen Chapman, Director of Hastings' Office of Student Services, rejected the application because CLS's bylaws excluded students on the bases of religion and sexual orientation. Hastings denied CLS's formal request for an exemption, making CLS the only student group ever denied RSO registration among more than 60 recognized groups at the school. Hastings stated, however, that CLS could still use Law School facilities for meetings and advertise on chalkboards and generally available campus bulletin boards. Refusing to alter its bylaws, CLS operated independently during the 2004--2005 academic year, holding weekly Bible studies, lectures, and social events, though the dissent later highlighted evidence that Hastings ignored or belatedly answered CLS's requests for a patio table and a guest-speaker room. The dispute unfolded against a broader national conflict between campus nondiscrimination mandates protecting LGBTQ students and religious organizations' claims of associational autonomy under the First and Fourteenth Amendments.

\section{Procedural History}

On October 22, 2004, CLS filed suit under 42 U.S.C. \S~1983 in the United States District Court for the Northern District of California against Hastings officers and administrators (initially naming Dean Mary Kay Kane, and later Acting Dean Leo P. Martinez), seeking declaratory and injunctive relief. CLS alleged that denying RSO status violated its First and Fourteenth Amendment rights to free speech, expressive association, and free exercise of religion. Hastings Outlaw, an RSO dedicated to combating sexual-orientation discrimination, intervened as a defendant.

On cross-motions for summary judgment, the District Court ruled for Hastings on all claims (\href{https://supreme.justia.com/cases/federal/us/561/661/}{Justia}). It held that Hastings' all-comers condition on access to a limited public forum was both reasonable and viewpoint neutral; that denying RSO status did not unconstitutionally burden expressive association because Hastings withheld subsidies rather than compelling CLS to admit anyone; and that the policy satisfied the Free Exercise Clause under \textit{Employment Division v. Smith}, 494 U.S. 872 (1990), as a neutral, generally applicable rule.

The United States Court of Appeals for the Ninth Circuit affirmed in a two-sentence unpublished memorandum opinion (\textit{Christian Legal Soc. Chapter of Univ. of Cal. v. Kane}, 319 Fed. Appx. 645 (9th Cir. 2009)), holding that because the parties stipulated that Hastings imposed an all-comers open-membership rule on all student groups, the recognition conditions were viewpoint neutral and reasonable under \textit{Truth v. Kent School District}, 542 F.3d 634, 649--50 (9th Cir. 2008).

The Supreme Court granted certiorari on December 7, 2009 (No. 08-1371), heard oral argument on April 19, 2010, and on June 28, 2010, affirmed the Ninth Circuit's ruling that the all-comers policy is constitutional, while remanding the case for the Ninth Circuit to consider CLS's pretext and selective-enforcement argument if, and to the extent, preserved (\href{https://www.courtlistener.com/opinion/150544/christian-legal-soc-chapter-of-univ-of-cal-hastings-college-of-law-v/}{CourtListener}).

\section{Judicial Votes}

The Supreme Court decided the case by a 5--4 vote:

\begin{itemize}
  \item \textbf{Majority (5 Justices):} Justice Ginsburg delivered the opinion of the Court, joined in full by Justices Stevens, Kennedy, Breyer, and Sotomayor.
  \item \textbf{Concurrences (2 separate opinions):}
  \begin{itemize}
    \item Justice Stevens filed a concurring opinion (joining the Court's opinion without reservation).
    \item Justice Kennedy filed a concurring opinion (joining the Court's opinion).
  \end{itemize}
  \item \textbf{Dissent (4 Justices):} Justice Alito filed a dissenting opinion, joined in full by Chief Justice Roberts and Justices Scalia and Thomas.
\end{itemize}

\section{Holding}

\textbf{Judgment:} Affirmed and remanded.

\textbf{Rule Announced:} When a public university opens a student-organization recognition program as a limited public forum, both free-speech and expressive-association challenges to membership conditions are analyzed together under limited-public-forum doctrine rather than strict scrutiny. A public university may constitutionally require every recognized student organization to accept ``all comers'' as members and leadership candidates, because such an open-access rule is reasonable in light of the educational purposes of the forum and is viewpoint neutral. Moreover, because the policy is neutral and generally applicable, the Free Exercise Clause does not entitle religious student groups to an exemption under \textit{Employment Division v. Smith}.

\textbf{Holding vs. Dicta:} The narrow holding applies solely to a genuine ``all-comers'' policy as stipulated by the parties; the Court expressly declined to decide the constitutionality of Hastings' written Nondiscrimination Policy standing alone (which prohibited discrimination on enumerated grounds such as religion and sexual orientation while arguably permitting ideological selectivity by secular groups). Consequently, Justice Stevens's defense of the written Nondiscrimination Policy in his concurrence, as well as the majority's speculation that student saboteurs are unlikely to hijack minority groups and Justice Kennedy's suggestion that proven hostile infiltration or pretextual enforcement would present a different case, constitute dicta or issues left open for remand and future litigation.

\section{Analysis of Opinions}

\subsection{Majority Opinion (Justice Ginsburg)}

Justice Ginsburg first addressed a threshold procedural issue: CLS urged the Court to evaluate Hastings' written Nondiscrimination Policy rather than the all-comers interpretation, arguing that the written text singled out religious and sexual-morality groups while allowing political clubs to exclude ideological opponents. The Court rejected this effort as an ``unseemly attempt to escape'' the December 2005 Joint Stipulation, emphasizing that formal factual stipulations are binding and conclusive judicial admissions (\href{https://supreme.justia.com/cases/federal/us/561/661/}{Justia}).

On the merits, the majority merged CLS's free-speech and expressive-association claims under limited-public-forum doctrine, declining to apply the strict scrutiny used in forced-inclusion cases like \textit{Roberts v. United States Jaycees}, 468 U.S. 609 (1984), and \textit{Boy Scouts of America v. Dale}, 530 U.S. 640 (2000). Ginsburg gave three reasons: (1) speech and expressive association are intertwined in this context, and it would be anomalous for a rule to survive limited-public-forum review as a speech restriction only to fail strict scrutiny as an associational restriction; (2) strict scrutiny would destroy a defining feature of limited public forums---the State's power to reserve the forum for specific classes of speakers, such as students only; and (3) Hastings was ``dangling the carrot of subsidy, not wielding the stick of prohibition,'' leaving CLS free to exclude anyone it wished outside the RSO program.

Applying limited-public-forum review, the Court distinguished \textit{Healy v. James}, 408 U.S. 169 (1972), \textit{Widmar v. Vincent}, 454 U.S. 263 (1981), and \textit{Rosenberger v. Rector and Visitors of Univ. of Va.}, 515 U.S. 819 (1995), because those universities targeted groups for disfavored treatment based on their viewpoints. Extending ``decent respect'' to educators' pedagogical judgments, Ginsburg found Hastings' all-comers policy reasonable on four independent grounds: (1) it ensures that leadership and educational opportunities funded by mandatory student-activity fees are open to every fee-paying student; (2) it relieves administrators of the ``daunting labor'' of policing whether belief-based exclusions mask prohibited status discrimination (noting that the Court has declined to separate homosexual status from homosexual conduct under \textit{Lawrence v. Texas}, 539 U.S. 558 (2003)); (3) it promotes tolerance, cooperation, and conflict-resolution skills across diverse viewpoints; and (4) it aligns with California's state-law antidiscrimination policy. Substantial alternative channels---including campus meeting rooms, chalkboards, bulletin boards, and electronic media---further minimized any burden on CLS.

Finally, the Court deemed the all-comers rule ``textbook viewpoint neutral.'' Even if it disproportionately burdened groups with unpopular campus views, under \textit{Ward v. Rock Against Racism}, 491 U.S. 781 (1989), a regulation justified without reference to expressive content does not become viewpoint-based merely because of differential incidental effects. In a brief footnote (n. 27), the Court dismissed CLS's Free Exercise Clause claim as foreclosed by \textit{Smith}, and in Part IV it remanded CLS's selective-enforcement and pretext claim for the Ninth Circuit to address if preserved.

\subsection{Concurring Opinion (Justice Stevens)}

Justice Stevens joined the Court's opinion in full but wrote separately to refute Justice Alito's contention that Hastings' written Nondiscrimination Policy would be ``plainly unconstitutional'' if reviewed on its face (\href{https://supreme.justia.com/cases/federal/us/561/661/}{Justia}). Stevens argued that even the written Nondiscrimination Policy is content and viewpoint neutral because it targets discriminatory conduct rather than belief or expression, applies equally to secular and religious acts of religious discrimination, and reflects a reasonable pedagogical judgment that invidious discrimination on grounds like race, religion, and sexual orientation is unacceptable in a university-sponsored program. Emphasizing that a university campus is ``a world apart from the public square,'' Stevens concluded that while a free society must tolerate exclusionary groups off campus, ``[i]t need not subsidize them, give them its official imprimatur, or grant them equal access to law school facilities.''

\subsection{Concurring Opinion (Justice Kennedy)}

Justice Kennedy, the author of \textit{Rosenberger}, wrote separately to explain why \textit{Rosenberger} did not control and to underscore the educational rationale supporting Hastings' policy (\href{https://supreme.justia.com/cases/federal/us/561/661/}{Justia}). Unlike the policy in \textit{Rosenberger}, Hastings' all-comers rule was not content-based in formulation or purpose; instead, it served the legitimate educational aim of encouraging law students to engage across differences in extracurricular life, noting that ``[a] vibrant dialogue is not possible if students wall themselves off from opposing points of view'' and that ``[t]he era of loyalty oaths is behind us.'' Crucially, however, Kennedy added a limiting caveat: CLS would have ``a substantial case on the merits'' if it could demonstrate that a school adopted or enforced an all-comers policy with intent to disadvantage a group for its views, or if hostile students actually used the policy to infiltrate a group and stifle its message.

\subsection{Dissenting Opinion (Justice Alito)}

Justice Alito, joined by Chief Justice Roberts and Justices Scalia and Thomas, delivered a sharp dissent accusing the majority of endorsing ``no freedom for expression that offends prevailing standards of political correctness in our country's institutions of higher learning'' (\href{https://www.law.cornell.edu/supct/html/08-1371.ZD.html}{Cornell LII}). Alito first contested the majority's factual framing, marshaling record evidence that Hastings actually denied CLS's application under its written Nondiscrimination Policy, admitted in its May 2005 Answer that secular political and cultural groups could require ideological adherence, and only unveiled the undocumented ``all-comers'' interpretation in a July 2005 deposition---before shifting again in its Supreme Court brief to an ``accept-some-comers'' policy permitting neutral conduct and skill rules. Alito also emphasized that Hastings had routinely registered groups (such as the Hastings Democratic Caucus, Silenced Right, and La Raza) whose bylaws restricted membership by viewpoint or identity, and that Hastings had stonewalled CLS's requests for a patio table and meeting room.

Doctrinally, Alito argued that \textit{Healy v. James} was directly controlling: in \textit{Healy}, the Court refused to defer to university administrators and held that denying official recognition and campus communication channels substantially burdened expressive association regardless of off-campus alternatives. Even under limited-public-forum doctrine, Alito contended that (1) the written Nondiscrimination Policy was patent viewpoint discrimination because it allowed secular groups to enforce ideological alignment while forbidding religious groups from doing the same; (2) the all-comers policy was unreasonable in light of the RSO forum's stipulated purpose of promoting viewpoint diversity \textit{among} student groups; (3) the record presented overwhelming evidence of pretext (shifting rationales, litigation-driven timing, lack of documentation, and prior nonenforcement); and (4) the all-comers rule exposed small, unpopular religious groups to hostile takeovers and ``marginalization,'' undermining the ``uninhibited, robust, and wide-open'' debate guaranteed by the First Amendment.

\section{Examples: Future Applications}

\begin{enumerate}
  \item \textbf{Same Side (Constitutional under \textit{Martinez}):} A public state university adopts a written, uniformly enforced all-comers policy for its student-activity-fee forum from the outset of the program, requiring all 120 registered clubs---including political parties, advocacy groups, and religious fellowships---to open voting membership and officer elections to every enrolled student while allowing neutral attendance and dues rules. A conservative political club refuses to remove its requirement that officers sign a statement supporting free-market principles and is denied RSO funding. Under \textit{Martinez}, the university prevails because the policy is genuinely applied across the board, reasonable in light of the forum's educational goals, and viewpoint neutral.

  \item \textbf{Same Side (Constitutional under \textit{Martinez}):} A public municipal parks department operates a subsidized community-center room-reservation and grant program open to local civic associations on the strict condition that any participating club allow any city resident to join and run for board office regardless of beliefs or background. An atheist civic club that restricts membership to nonbelievers is denied subsidized grants and priority booking, though it may still rent rooms on standard public terms. Because the city is dangling a subsidy within a limited public forum rather than prohibiting private association, and enforces a neutral all-comers condition uniformly, the condition survives under \textit{Martinez}.

  \item \textbf{Opposite Side (Unconstitutional despite \textit{Martinez}):} A public university adopts a ``Human Rights Policy'' prohibiting student groups from selecting leaders based on religion or sexual orientation, and deregisters a Christian student group (such as Business Leaders in Christ or InterVarsity Christian Fellowship) for requiring its student leaders to affirm its faith, while simultaneously allowing secular political, ideological, and identity-based clubs to require leaders to support their missions and granting exemptions to fraternities and sports clubs. As the Eighth Circuit held in \href{https://becketfund.org/case/intervarsity-christian-fellowship-v-university-iowa/}{\textit{InterVarsity Christian Fellowship v. University of Iowa}} (2021), this selective enforcement falls outside \textit{Martinez}'s all-comers safe harbor and constitutes unconstitutional viewpoint discrimination.

  \item \textbf{Opposite Side (Unconstitutional despite \textit{Martinez}):} A state legislature enacts a general public-accommodations statute requiring every private, off-campus religious fellowship and Bible-study group operating in the state to admit any resident as a voting member and officer regardless of faith or conduct, backed by civil fines. Unlike Hastings' RSO program, this law wields ``the stick of prohibition'' rather than ``the carrot of subsidy,'' triggering strict scrutiny under \textit{Boy Scouts of America v. Dale} and \textit{Roberts v. United States Jaycees}, and is plainly unconstitutional under the First Amendment.

  \item \textbf{Fence-Sitter (Hard Case):} A public law school genuinely adopts and uniformly enforces a written all-comers policy for all RSOs. A small, seven-member Orthodox Jewish student group complies in order to retain campus meeting rooms and funding. Soon thereafter, fifteen ideologically opposed students pay the club's nominal dues, attend two meetings, vote themselves into all officer positions, and convert the group into an anti-Zionist advocacy club. The original founders sue, invoking Justice Kennedy's concurrence in \textit{Martinez}, which stated that a group ``would have a substantial case on the merits'' if an all-comers policy were ``used to infiltrate the group or challenge its leadership in order to stifle its views.'' Because the majority opinion dismissed hostile takeovers as hypothetical and did not specify the doctrinal test for as-applied infiltration challenges, lower courts would face a genuine doctrinal impasse between deferring to a facially neutral all-comers rule and protecting minority speakers from a majoritarian heckler's veto.
\end{enumerate}

\section{Critique}

\textit{Christian Legal Society v. Martinez} has drawn sustained criticism across the ideological spectrum, with scholars highlighting both its doctrinal vulnerabilities and its competing constitutional premises (\href{https://fedsoc.org/case/christian-legal-society-v-martinez}{Federalist Society}; \href{https://lawreview.law.pitt.edu/ojs/lawreview/article/download/1172/737/2196}{Pittsburgh Law Review}).

\textbf{Logical and Doctrinal Weaknesses:} First, critics note that the Court's heavy reliance on the parties' Joint Stipulation allowed it to decide an artificial, stylized case while ducking the far more common real-world controversy: university nondiscrimination policies that bar religious and sexual-orientation criteria while permitting secular ideological screening. Second, merging expressive association into limited-public-forum review creates an internal tension: a limited public forum for student clubs exists to foster diverse viewpoints \textit{among} groups, yet an all-comers rule can logically undermine that very purpose by preventing minority groups from maintaining ideological coherence. Third, the majority's ``carrot of subsidy'' versus ``stick of prohibition'' distinction sits uneasily alongside the unconstitutional conditions doctrine and \textit{Rosenberger}, which treated student-activity-fee forums as vital channels of campus speech rather than discretionary government favors. Finally, Justice Kennedy's concurrence left an unresolved analytical gap: if an all-comers policy is facially reasonable and viewpoint neutral, it is unclear why private students' lawful use of that policy to outvote a small club would render the \textit{university's} policy unconstitutional state action.

\textbf{Originalist and Textualist Critique:} Originalist, textualist, and religious-liberty scholars (including John Inazu in \textit{Liberty's Refuge: The Forgotten Freedom of Assembly} and Michael W. McConnell; see \href{https://digitalcommons.pepperdine.edu/cgi/viewcontent.cgi?article=2355&context=plr}{Pepperdine Law Review}) fault \textit{Martinez} for subordinating the First Amendment's textually enumerated right ``peaceably to assemble'' and its Free Exercise guarantee to modern, judge-made forum tiers. From this perspective, the Founders understood religious associations and voluntary assemblies as independent, self-defining communities whose autonomy necessarily entails the power to define membership by shared creed. Treating a religious group's faith-based leadership requirement as mere ``conduct'' rather than core protected exercise, and deferring to university administrators' pedagogical preferences, reverses the First Amendment's baseline skepticism of state orthodoxy. Subsequent Roberts Court decisions strengthening Free Exercise protections against selective exemptions (such as \textit{Fulton v. City of Philadelphia}) and Eighth Circuit rulings like \href{https://becketfund.org/case/intervarsity-christian-fellowship-v-university-iowa/}{\textit{InterVarsity Christian Fellowship v. University of Iowa}} have increasingly cabined \textit{Martinez} to strictly verified, uniformly applied all-comers policies.

\textbf{Progressive and Egalitarian Perspectives:} Progressive constitutional scholars generally defend \textit{Martinez}'s outcome as an essential affirmation of civic equality and democratic pluralism on public campuses, arguing that public universities need not compel LGBTQ or minority-faith students to subsidize groups that exclude them from full membership and leadership (\href{https://brooklynworks.brooklaw.edu/cgi/viewcontent.cgi?article=1051&context=blr}{Brooklyn Law Review}). Yet even some progressive and civil-liberties commentators critique the majority's reasoning from the left. First, endorsing a sweeping ``all-comers'' mandate rather than Justice Stevens's narrower defense of targeted antidiscrimination rules threatens the associational integrity of progressive affinity and identity groups (such as Black, feminist, or LGBTQ student organizations), which may equally fear dilution or disruption if barred from requiring members to share their core mission. Second, granting broad judicial deference to university administrators in campus speech disputes risks weakening \textit{Healy v. James} and empowering hostile state legislatures or university boards to restrict disfavored progressive student movements through facially neutral administrative rules.

\section{Key Quotations}

\begin{quote}
The First Amendment shields CLS against state prohibition of the organization's expressive activity, however exclusionary that activity may be. But CLS enjoys no constitutional right to state subvention of its selectivity.
\begin{flushright}--- Justice Ginsburg, Opinion of the Court (\href{https://supreme.justia.com/cases/federal/us/561/661/}{561 U.S. at 669})\end{flushright}
\end{quote}

\begin{quote}
Application of the less-restrictive limited-public-forum analysis better accounts for the fact that Hastings, through its RSO program, is dangling the carrot of subsidy, not wielding the stick of prohibition.
\begin{flushright}--- Justice Ginsburg, Opinion of the Court (\href{https://supreme.justia.com/cases/federal/us/561/661/}{561 U.S. 661})\end{flushright}
\end{quote}

\begin{quote}
It is, after all, hard to imagine a more viewpoint-neutral policy than one requiring all student groups to accept all comers.
\begin{flushright}--- Justice Ginsburg, Opinion of the Court (\href{https://supreme.justia.com/cases/federal/us/561/661/}{561 U.S. 661})\end{flushright}
\end{quote}

\begin{quote}
A vibrant dialogue is not possible if students wall themselves off from opposing points of view.
\begin{flushright}--- Justice Kennedy, Concurring (\href{https://supreme.justia.com/cases/federal/us/561/661/}{561 U.S. 661})\end{flushright}
\end{quote}

\begin{quote}
The proudest boast of our free speech jurisprudence is that we protect the freedom to express ``the thought that we hate.'' \textit{United States v. Schwimmer}, 279 U. S. 644, 654--655 (1929) (Holmes, J., dissenting). Today's decision rests on a very different principle: no freedom for expression that offends prevailing standards of political correctness in our country's institutions of higher learning.
\begin{flushright}--- Justice Alito, Dissenting (\href{https://www.law.cornell.edu/supct/html/08-1371.ZD.html}{561 U.S. 661})\end{flushright}
\end{quote}

\end{document}
